\section{Continuous Assessment 1 --- Set B}
\label{sec:ca1_set_b}

\LPUCSHeader{Continuous Assessment 1 (Objective MCQ)}{B}{50 Minutes}{30 Marks}{CO1 \& CO2 (Unit 1: Computer System \& Unit 2: Operating System)}

\begin{center}
\sffamily\textbf{\color{AlertRed}Instructions: All 30 questions are compulsory. Each question carries 1 mark. A penalty of 0.25 marks is deducted for each incorrect response.}
\end{center}

\vspace{3mm}

% ------------------------------------------------------------------------------
% UNIT 1: COMPUTER SYSTEM (Q1 TO Q15)
% ------------------------------------------------------------------------------
\mcqitem{1}{The fundamental sequence of operations performed by the CPU to process a single machine language instruction is:}{Fetch $\to$ Decode $\to$ Execute $\to$ Writeback}{Decode $\to$ Fetch $\to$ Compile $\to$ Execute}{Link $\to$ Assemble $\to$ Execute $\to$ Store}{Interrupt $\to$ Fetch $\to$ Decode $\to$ Link}{1}

\mcqitem{2}{The central processor register that holds the specific machine instruction currently being decoded and executed by the control unit is the:}{Instruction Register (IR)}{Program Counter (PC)}{Memory Address Register (MAR)}{Stack Pointer (SP)}{1}

\mcqitem{3}{In modern multi-core processor architectures, Level 1 (L1) cache is typically:}{Split into separate L1 Instruction (L1i) and L1 Data (L1d) caches dedicated to each CPU core}{Shared globally across all cores on the motherboard}{Located outside the silicon chip on DIMM sockets}{Slower than Level 3 (L3) cache}{1}

\mcqitem{4}{Under a Write-Back cache policy, when the CPU modifies a cached memory location, the updated data is written to main memory:}{Only when the modified cache block is evicted from the cache}{Immediately on every single write operation}{Only upon system shutdown}{Never, as data remains exclusively in registers}{1}

\mcqitem{5}{Solid-State Drives (SSDs) employ Wear Leveling algorithms in their Flash Translation Layer (FTL) to:}{Distribute write and erase cycles evenly across all NAND blocks to prevent premature drive failure}{Speed up mechanical disk rotation}{Compress files automatically on write}{Double the capacity of the RAM}{1}

\mcqitem{6}{Which type of NAND flash memory stores exactly 1 bit per memory cell, providing the highest endurance and read/write speeds?}{SLC (Single-Level Cell)}{MLC (Multi-Level Cell)}{TLC (Triple-Level Cell)}{QLC (Quad-Level Cell)}{1}

\mcqitem{7}{In a RAID 1 array comprising two $2\text{ TB}$ hard drives, the total usable storage capacity available to the operating system is:}{$2\text{ TB}$ (50\% storage efficiency due to exact data mirroring)}{$4\text{ TB}$}{$1\text{ TB}$}{$3\text{ TB}$}{1}

\mcqitem{8}{Which RAID level utilizes two independent distributed parity blocks ($P$ and $Q$) per stripe to tolerate the simultaneous failure of any two disks?}{RAID 6}{RAID 5}{RAID 0}{RAID 1}{1}

\mcqitem{9}{Instruction pipelining in modern CPU design enhances processor performance primarily by:}{Overlapping the execution phases of consecutive instructions to increase overall instruction throughput}{Increasing the clock frequency without generating heat}{Eliminating all branch statements from code}{Removing the need for a cache hierarchy}{1}

\mcqitem{10}{A hardware interrupt generated by an external peripheral (e.g., keyboard keystroke or network packet arrival) causes the CPU to:}{Suspend current execution and invoke the appropriate Interrupt Service Routine (ISR)}{Immediately reboot the motherboard}{Crash the active operating system}{Disable the memory bus permanently}{1}

\mcqitem{11}{If a CPU operates at a clock frequency of $4.0\text{ GHz}$, the duration of a single clock period (cycle time $T$) is:}{$0.25\text{ nanoseconds}$ ($250\text{ ps}$)}{$4.0\text{ nanoseconds}$}{$0.25\text{ microseconds}$}{$2.5\text{ milliseconds}$}{1}

\mcqitem{12}{A PCIe 4.0 slot provides a unidirectional data bandwidth of approximately $2\text{ GB/s}$ per lane. What is the theoretical unidirectional bandwidth of a PCIe 4.0 x16 graphics card slot?}{$\approx 32\text{ GB/s}$}{$\approx 16\text{ GB/s}$}{$\approx 64\text{ GB/s}$}{$\approx 8\text{ GB/s}$}{1}

\mcqitem{13}{Which of the following computer storage media is fundamentally volatile, losing all stored data when electrical power is disconnected?}{DRAM (Dynamic Random-Access Memory)}{NAND Flash Drive}{Optical Blu-ray Disc}{Magnetic Hard Disk Drive}{1}

\mcqitem{14}{The hardware interface mechanism that allows the CPU to access memory and I/O devices using the same address space and identical machine instructions is called:}{Memory-Mapped I/O (MMIO)}{Port-Mapped I/O}{Isolated I/O}{Direct Cache Invalidation}{1}

\mcqitem{15}{In symmetric multiprocessing (SMP) computer systems, cache coherence protocols (such as MESI) are required to ensure that:}{Multiple CPU cores reading from their private caches see a consistent, identical view of shared main memory}{RAM data is refreshed every 64 milliseconds}{The GPU executes 3D graphics rendering}{Virtual memory pages are never swapped to disk}{1}

% ------------------------------------------------------------------------------
% UNIT 2: OPERATING SYSTEM (Q16 TO Q30)
% ------------------------------------------------------------------------------
\mcqitem{16}{The Linux operating system kernel is architecturally classified as a:}{Monolithic kernel with dynamic Loadable Kernel Module (LKM) capabilities}{Pure microkernel architecture}{Exokernel running exclusively in user space}{Hypervisor without device drivers}{1}

\mcqitem{17}{Dual-mode CPU execution (User Mode vs. Kernel Mode) provides operating system security by:}{Preventing unauthorized user applications from executing privileged hardware instructions or modifying kernel memory}{Allowing all user programs unrestricted direct access to device drivers}{Doubling the physical clock frequency of the processor}{Translating C code into Python automatically}{1}

\mcqitem{18}{Threads belonging to the same parent process share which of the following system resources?}{Code segment, data segment, and open file descriptors}{Program Counter and private CPU registers}{Execution stack and stack pointer}{Process identifier (PID)}{1}

\mcqitem{19}{Which of the following conditions is NOT one of the four necessary Coffman conditions required for a system deadlock to occur?}{Preemptive resource deallocation}{Mutual Exclusion}{Hold and Wait}{Circular Wait}{1}

\mcqitem{20}{In CPU process scheduling, the metric defined as the total time elapsed from process submission until its complete termination is the:}{Turnaround Time}{Waiting Time}{Response Time}{Throughput}{1}

\mcqitem{21}{The Convoy Effect, where several short CPU-burst processes are delayed while waiting for a single long CPU-intensive process to release the processor, is a major limitation of:}{First-Come, First-Served (FCFS) Scheduling}{Round Robin (RR) Scheduling}{Shortest Remaining Time First (SRTF)}{Multilevel Feedback Queue}{1}

\mcqitem{22}{Demand Paging in virtual memory architectures refers to the policy where:}{A virtual memory page is loaded into physical RAM only when an active process attempts to reference it}{All program pages are pre-loaded into RAM before the program starts}{Pages are permanently locked in high-speed registers}{Memory frames are deleted upon page faults}{1}

\mcqitem{23}{The Translation Lookaside Buffer (TLB) is a specialized hardware component used in paging systems to:}{Cache recent virtual-to-physical address translations to accelerate memory access}{Store CPU registers during an interrupt}{Hold uncompressed audio frames}{Schedule real-time threads}{1}

\mcqitem{24}{Internal fragmentation in memory allocation occurs when:}{The memory block allocated to a process is slightly larger than the requested size, leaving unused space inside the partition}{Free memory blocks are scattered across RAM but are individually too small to satisfy a request}{The disk drive runs out of sectors}{Two processes access the same shared variable}{1}

\mcqitem{25}{A Soft Real-Time system is exemplified by which of the following applications?}{Live multimedia video streaming where an occasional dropped packet degrades quality but does not destroy the system}{An aircraft fly-by-wire flight control system}{A nuclear reactor control rod trip mechanism}{An automotive anti-lock braking system (ABS)}{1}

\mcqitem{26}{Which software compilation utility translates high-level source code into assembly language or machine-specific object code while performing optimizations?}{Compiler}{Linker}{Loader}{Device Driver}{1}

\mcqitem{27}{Just-In-Time (JIT) compilation, used by environments such as the Java Virtual Machine (JVM) and V8 engine, improves execution performance by:}{Compiling frequently executed bytecode into native machine code dynamically during runtime}{Translating source code to PDF documents}{Executing assembly mnemonics on an external GPU}{Eliminating all variables at compile time}{1}

\mcqitem{28}{In software development, Dynamic Linking (shared libraries like \texttt{.so} in Linux or \texttt{.dll} in Windows) offers the benefit of:}{Allowing multiple concurrent processes to share a single copy of library code in physical RAM}{Producing self-contained executable binaries that require no runtime dependencies}{Eliminating the need for an operating system loader}{Preventing any runtime updates to shared libraries}{1}

\mcqitem{29}{In the standard virtual address space layout of a running C process, dynamically allocated memory requested via \texttt{malloc()} is allocated on the:}{Heap segment (which expands upwards towards higher memory)}{Stack segment (which holds local function variables)}{Text segment (which stores read-only machine instructions)}{BSS segment (which stores uninitialized static globals)}{1}

\mcqitem{30}{Stack memory allocation differs fundamentally from Heap memory allocation because stack memory:}{Operates in a strict Last-In, First-Out (LIFO) order with automatic allocation and deallocation on function entry and return}{Requires manual deallocation via \texttt{free()} to avoid memory leaks}{Supports arbitrary dynamic resizing during runtime without size limits}{Is stored on secondary hard disk storage rather than physical RAM}{1}

\vspace{5mm}
\hrule
\vspace{5mm}

% ------------------------------------------------------------------------------
% ANSWER KEY & EXPLANATIONS
% ------------------------------------------------------------------------------
\subsection*{Answer Key (CA-1 Set B)}

\begin{center}
\begin{tabular}{|c|c||c|c||c|c||c|c||c|c||c|c|}
\hline
\textbf{Q} & \textbf{Ans} & \textbf{Q} & \textbf{Ans} & \textbf{Q} & \textbf{Ans} & \textbf{Q} & \textbf{Ans} & \textbf{Q} & \textbf{Ans} & \textbf{Q} & \textbf{Ans} \\
\hline
1 & \textbf{A} & 6 & \textbf{A} & 11 & \textbf{A} & 16 & \textbf{A} & 21 & \textbf{A} & 26 & \textbf{A} \\
2 & \textbf{A} & 7 & \textbf{A} & 12 & \textbf{A} & 17 & \textbf{A} & 22 & \textbf{A} & 27 & \textbf{A} \\
3 & \textbf{A} & 8 & \textbf{A} & 13 & \textbf{A} & 18 & \textbf{A} & 23 & \textbf{A} & 28 & \textbf{A} \\
4 & \textbf{A} & 9 & \textbf{A} & 14 & \textbf{A} & 19 & \textbf{A} & 24 & \textbf{A} & 29 & \textbf{A} \\
5 & \textbf{A} & 10 & \textbf{A} & 15 & \textbf{A} & 20 & \textbf{A} & 25 & \textbf{A} & 30 & \textbf{A} \\
\hline
\end{tabular}
\end{center}

\vspace{3mm}
\subsection*{Detailed Technical Solutions}

\begin{solutionbox}[Key Architectural \& Systems Explanations]
\textbf{Q.11 (Ans: A)} The clock period $T = \frac{1}{f} = \frac{1}{4.0 \times 10^9\text{ s}^{-1}} = 0.25 \times 10^{-9}\text{ s} = 0.25\text{ ns} = 250\text{ ps}$.\\[2mm]
\textbf{Q.12 (Ans: A)} PCIe 4.0 transfers at $16\text{ GT/s}$ with a 128b/130b encoding scheme, delivering $\approx 1.969\text{ GB/s}$ ($\approx 2\text{ GB/s}$) per lane. Across 16 lanes ($x16$), the unidirectional bandwidth is $16 \times 2\text{ GB/s} = 32\text{ GB/s}$ (bidirectional total $\approx 64\text{ GB/s}$).\\[2mm]
\textbf{Q.18 (Ans: A)} Threads created within the same process share the process address space (code/text, data, heap) as well as operating system resources like open file descriptors and signals. However, each thread maintains its own thread ID, registers (including PC), and private call stack.\\[2mm]
\textbf{Q.19 (Ans: A)} The 4 necessary Coffman conditions for deadlock are: (1) Mutual Exclusion, (2) Hold and Wait, (3) \textbf{No Preemption} (resources cannot be forcibly taken from a process), and (4) Circular Wait. Therefore, "Preemptive resource deallocation" prevents deadlock rather than causing it.\\[2mm]
\textbf{Q.23 (Ans: A)} Without a TLB, every virtual memory reference requires accessing the page table in physical RAM before fetching the actual target data (2 memory accesses). The TLB provides an associative lookup cache that resolves virtual page numbers to physical frames in $\sim 1\text{ ns}$, yielding effective memory access times close to that of physical RAM.\\[2mm]
\textbf{Q.29 (Ans: A)} In standard Unix/Linux memory layout, the lowest addresses contain the read-only Text (code) segment, followed by Initialized Data, Uninitialized Data (BSS), and then the Heap (which grows upward towards high addresses). The Stack starts at high memory and grows downwards.
\end{solutionbox}

\newpage
